%------------------------------------------------------------------------------%
\documentclass[fleqn,utf8,aspectratio=169,12pt]{beamer}

\usepackage{integracao_e_series}

\title{O Método (Secreto) da Matemática}

%------------------------------------------------------------------------------%
\begin{document}

\addtitle

%------------------------------------------------------------------------------%
\section{Método Matemático}

%------------------------------------------------------------------------------%
\begin{frame}{O Método Matemático}
  Não basta encontrar a solução correta, \\[2mm]
  precisamos da solução \emph{comprovadamente correta}
\end{frame}

%------------------------------------------------------------------------------%
\begin{frame}{O Método Matemático}
  Uma aproximação não é uma solução
\end{frame}

%------------------------------------------------------------------------------%
\begin{frame}{O Método Matemático}
  Não podemos usar o método empírico
\end{frame}

%------------------------------------------------------------------------------%
\begin{frame}{O Método Matemático}
  Toda afirmação precisa ser \emph{comprovadamente verdadeira}
\end{frame}

%------------------------------------------------------------------------------%
\begin{frame}{O Método Matemático}
  Detalhes são extremamente importantes

  \pause
  Notações similares podem significar coisas muito diferentes
\end{frame}

%------------------------------------------------------------------------------%
\begin{frame}{O Método Matemático}
  \begin{enumerate}[<+->]
    \item Axiomas
    \begin{itemize}
      \item Definem uma teoria
      \item Combinamos que é verdade
    \end{itemize}
    ~

    \item Definições
    \begin{itemize}
      \item Definem conceitos em uma teoria
      \item Combinamos que é verdade
    \end{itemize}
    ~

    \item Proposições, Lemas, Teoremas, Corolários
    \begin{itemize}
      \item Demonstramos que é verdade
    \end{itemize}
  \end{enumerate}
\end{frame}

%------------------------------------------------------------------------------%
\begin{frame}{O Método Matemático}
  Nem sempre o \textit{``obvio''} é verdadeiro
\end{frame}

%------------------------------------------------------------------------------%
\begin{frame}{O Método Matemático}
  Cada passagem em um desenvolvimento deve \\
  \emph{obrigatoriamente}
  estar ancorada em um resultado demonstrado

  \pause
  \bigskip
  Por isso eles são numerados
\end{frame}

%------------------------------------------------------------------------------%
\section{Resolvendo Exercícios em uma Prova}

%------------------------------------------------------------------------------%
\begin{frame}{Na Prova}
  \begin{itemize}[<+->]
    \item A resposta correta tem pouca importância
    \item Sua capacidade de manipular corretamente os conceitos é o que importa
    \item Você não quer depender
    \begin{itemize}
      \item da boa vontade do corretor
      \item da capacidade do corretor em decifrar o que você escreveu
      \item da capacidade do corretor em adivinhar o que você pensou
    \end{itemize}
  \end{itemize}
\end{frame}

%------------------------------------------------------------------------------%
\begin{frame}{Organização do Prova}
  \begin{itemize}[<+->]
    \item Respeite as margens do papel (Todas as Quatro!)
    \item Use uma ``fonte'' adequada, cada coisa tem seu símbolo \emph{único}
    \item Não imite o quadro do professor ou as imagens de um vídeo
    \item Em português escrevemos da esquerda para a direita e de cima para baixo
    \begin{itemize}
      \item siga essa regra
      \item não use setas para indicar o sentido de leitura
      \item a informação não pode vir do futuro
    \end{itemize}
  \end{itemize}
\end{frame}

%------------------------------------------------------------------------------%
\begin{frame}{Notação}
  \begin{itemize}[<+->]
    \item Use a notação correta
    \item Não misture os símbolos \; $=$,\; $\to$,\; $\Rightarrow$
    \begin{itemize}
      \item \emph{$=$}           \tabto{8mm} indica igualdade, ambos os lados são exatamente a mesma coisa
      \item \emph{$\to$}         \tabto{8mm} é usado somente na notação de limite e na definição de funções
      \item \emph{$\Rightarrow$} \tabto{8mm} indica implicação lógica \\[3mm]
        \onslide<+->{
        ( A terra é plana ) $\quad\Rightarrow\quad$ ( Belo Horizonte tem praias maravilhosas ) \\[3mm]}
        \onslide<+->{
        Afirmação \emph{verdadeira}, pois falso implica em falso}
    \end{itemize}
  \end{itemize}
\end{frame}

%------------------------------------------------------------------------------%
\begin{frame}{Exemplo do uso de $\Rightarrow$}
Afirmação verdadeira para $x\in\R$
\[
  x > 7 \quad\Rightarrow\quad x > 3
\]

\pause
Substituindo valores
\pause
\begin{ceqn}
\[
\begin{array}{>{\quad}c<{\quad}>{\quad}c<{\quad}>{\quad}c<{\quad}} \hline
  x        & x > 7             & x > 3             \\ \hline \pause
  9        & \text{verdadeiro} & \text{verdadeiro} \\ \pause
  5        & \text{falso}      & \text{verdadeiro} \\ \pause
  1        & \text{falso}      & \text{falso}      \\ \pause
  \nexists & \text{falso}      & \text{verdade}    \\ \hline
\end{array}
\]
\end{ceqn}
\end{frame}

%------------------------------------------------------------------------------%
\begin{frame}{Evite Erros Comuns}
\begin{columns}[t]
\column{0.3\linewidth}
  \begin{itemize}[<+->]
    \setlength\itemsep{1.2em}
    \item Não existe $\times-$
    \item \(f(x) \neq f(a)\)
    \item \(\frac{a+b}{c+b} \neq \frac{a}{c}+\frac{b}{d}\)
    \item \((a+b)^2 \neq a^2 + b^2\)
    \item \(\sqrt{4} = 2 \neq \pm 2 \)
    \item \(\sqrt{x^2} \neq x\)
  \end{itemize}
\column{0.7\linewidth}
  \begin{itemize}[<+->]
    \setlength\itemsep{1em}
    \item Não declarar o que está calculando
    \item Usar informação que só será apresentada no futuro
  \end{itemize}
\end{columns}
\end{frame}

%------------------------------------------------------------------------------%
\begin{frame}{Não Pule Passagens}
  \begin{itemize}[<+->]
    \item Escreva todas as passagens
    \item O solicitado no enunciado deve estar escrito explicitamente na resolução
    \item Evite repetir expressões grandes, atribua um nome para elas
  \end{itemize}
\end{frame}

%------------------------------------------------------------------------------%
\begin{frame}{Repetições}

Prefira sequências de igualdades a sequências de equações

\bigskip
\begin{columns}
\column{0.3\linewidth}
  \pause
  \quad Evite
  \vspace{-0.4\baselineskip}
  \begin{align*}
    f(x) & = 2 + 2 + 2 + 2 \\[-1mm]
    f(x) & = 4 + 2 + 2     \\[-1mm]
    f(x) & = 6 + 2         \\[-1mm]
    f(x) & = 8
  \end{align*}
\column{0.6\linewidth}
  \pause
  \quad Prefira
  \vspace{-0.4\baselineskip}
  \begin{align*}
    f(x) & = 2 + 2 + 2 + 2 \\[-1mm]
         & = 4 + 2 + 2     \\[-1mm]
         & = 6 + 2         \\[-1mm]
         & = 8
  \end{align*}
\end{columns}

\pause
\bigskip
Forma alternativa (menos legível)
\[
  f(x)
  = 2 + 2 + 2 + 2
  = 4 + 2 + 2
  = 6 + 2
  = 8
\]
\end{frame}

%------------------------------------------------------------------------------%
\begin{frame}{Demonstrando Igualdades}
  \onslide<+->{Começamos de um termo e o manipulamos até o outro}

  \bigskip
  \onslide<+->{Mostrar que \(a^2-b^2 = (a+b)(a-b)\)}

  \begin{align*}
    \onslide<+->{(a+b)(a-b)}
    \onslide<+->{& = a(a-b) + b(a-b) \\}
    \onslide<+->{& = a^2 - ab + ba - b^2 \\}
    \onslide<+->{& = a^2 - b^2}
  \end{align*}
\end{frame}

%------------------------------------------------------------------------------%
\begin{frame}{Demonstrando Igualdades}
  \onslide<+->{\emph{Não manipular os dois lados ao mesmo tempo}}

  \bigskip
  \onslide<+->{Mostrar que \(3=4\)}

  \begin{align*}
    \onslide<+->{3          & = 4          \\}
    \onslide<+->{3 \times 0 & = 4 \times 0 \\}
    \onslide<+->{0          & = 0}
  \end{align*}

  \onslide<+->{Se o erro não está claro para você, não manipule os dois lados}
\end{frame}

%------------------------------------------------------------------------------%
\section{Lista Mínima}

%------------------------------------------------------------------------------%
\begin{frame}{Lista Mínima}
  Ler o texto
  ``\href{https://drive.google.com/file/d/1iUuy68VWqKabrV7JxTwlyOgr59cA9G2r/view?usp=sharing}%
  {Como Resolver Exercícios de Matemática}''

  ~

  Ler a Apostila
  ``\href{https://material-didatico.github.io/pages/infinito/}%
  {Quem disse que $0,999\dots=1$}''
\end{frame}

%------------------------------------------------------------------------------%
\end{document}
%------------------------------------------------------------------------------%
